
\documentclass[a4paper,11pt]{article}
\usepackage[a4paper, margin=8em]{geometry}

% usa i pacchetti per la scrittura in italiano
\usepackage[french,italian]{babel}
\usepackage[T1]{fontenc}
\usepackage[utf8]{inputenc}
\frenchspacing 

% usa i pacchetti per la formattazione matematica
\usepackage{amsmath, amssymb, amsthm, amsfonts}

% usa altri pacchetti
\usepackage{gensymb}
\usepackage{hyperref}
\usepackage{standalone}

% imposta il titolo
\title{Appunti Comunicazioni Numeriche}
\author{Luca Seggiani}
\date{2026}

% imposta lo stile
% usa helvetica
\usepackage[scaled]{helvet}
% usa palatino
\usepackage{palatino}
% usa un font monospazio guardabile
\usepackage{lmodern}

\renewcommand{\rmdefault}{ppl}
\renewcommand{\sfdefault}{phv}
\renewcommand{\ttdefault}{lmtt}

% disponi il titolo
\makeatletter
\renewcommand{\maketitle} {
	\begin{center} 
		\begin{minipage}[t]{.8\textwidth}
			\textsf{\huge\bfseries \@title} 
		\end{minipage}%
		\begin{minipage}[t]{.2\textwidth}
			\raggedleft \vspace{-1.65em}
			\textsf{\small \@author} \vfill
			\textsf{\small \@date}
		\end{minipage}
		\par
	\end{center}

	\thispagestyle{empty}
	\pagestyle{fancy}
}
\makeatother

% disponi teoremi
\usepackage{tcolorbox}
\newtcolorbox[auto counter, number within=section]{theorem}[2][]{%
	colback=blue!10, 
	colframe=blue!40!black, 
	sharp corners=northwest,
	fonttitle=\sffamily\bfseries, 
	title=~\thetcbcounter: #2, 
	#1
}

% disponi definizioni
\newtcolorbox[auto counter, number within=section]{definition}[2][]{%
	colback=red!10,
	colframe=red!40!black,
	sharp corners=northwest,
	fonttitle=\sffamily\bfseries,
	title=~\thetcbcounter: #2,
	#1
}

% U.D.M
\newcommand{\amp}{\ensuremath{\mathrm{A}}}
\newcommand{\volt}{\ensuremath{\mathrm{V}}}
\newcommand{\meter}{\ensuremath{\mathrm{m}}}
\newcommand{\second}{\ensuremath{\mathrm{s}}}
\newcommand{\farad}{\ensuremath{\mathrm{F}}}
\newcommand{\henry}{\ensuremath{\mathrm{H}}}
\newcommand{\siemens}{\ensuremath{\mathrm{S}}}

% circuiti
\usepackage{circuitikz}
\usetikzlibrary{babel}

% disegni
\usepackage{pgfplots}
\pgfplotsset{width=10cm,compat=1.9}

% disponi codice
\usepackage{listings}
\usepackage[table]{xcolor}

\lstdefinestyle{codestyle}{
		backgroundcolor=\color{black!5}, 
		commentstyle=\color{codegreen},
		keywordstyle=\bfseries\color{magenta},
		numberstyle=\sffamily\tiny\color{black!60},
		stringstyle=\color{green!50!black},
		basicstyle=\ttfamily\footnotesize,
		breakatwhitespace=false,         
		breaklines=true,                 
		captionpos=b,                    
		keepspaces=true,                 
		numbers=left,                    
		numbersep=5pt,                  
		showspaces=false,                
		showstringspaces=false,
		showtabs=false,                  
		tabsize=2
}

\lstdefinestyle{shellstyle}{
		backgroundcolor=\color{black!5}, 
		basicstyle=\ttfamily\footnotesize\color{black}, 
		commentstyle=\color{black}, 
		keywordstyle=\color{black},
		numberstyle=\color{black!5},
		stringstyle=\color{black}, 
		showspaces=false,
		showstringspaces=false, 
		showtabs=false, 
		tabsize=2, 
		numbers=none, 
		breaklines=true
}

\lstdefinelanguage{javascript}{
	keywords={typeof, new, true, false, catch, function, return, null, catch, switch, var, if, in, while, do, else, case, break},
	keywordstyle=\color{blue}\bfseries,
	ndkeywords={class, export, boolean, throw, implements, import, this},
	ndkeywordstyle=\color{darkgray}\bfseries,
	identifierstyle=\color{black},
	sensitive=false,
	comment=[l]{//},
	morecomment=[s]{/*}{*/},
	commentstyle=\color{purple}\ttfamily,
	stringstyle=\color{red}\ttfamily,
	morestring=[b]',
	morestring=[b]"
}

% disponi sezioni
\usepackage{titlesec}

\titleformat{\section}
	{\sffamily\Large\bfseries} 
	{\thesection}{1em}{} 
\titleformat{\subsection}
	{\sffamily\large\bfseries}   
	{\thesubsection}{1em}{} 
\titleformat{\subsubsection}
	{\sffamily\normalsize\bfseries} 
	{\thesubsubsection}{1em}{}

% disponi alberi
\usepackage{forest}

\forestset{
	rectstyle/.style={
		for tree={rectangle,draw,font=\large\sffamily}
	},
	roundstyle/.style={
		for tree={circle,draw,font=\large}
	}
}

% disponi algoritmi
\usepackage{algorithm}
\usepackage{algorithmic}
\makeatletter
\renewcommand{\ALG@name}{Algoritmo}
\makeatother

% disponi numeri di pagina
\usepackage{fancyhdr}
\fancyhf{} 
\fancyfoot[L]{\sffamily{\thepage}}

\makeatletter
\fancyhead[L]{\raisebox{1ex}[0pt][0pt]{\sffamily{\@title \ \@date}}} 
\fancyhead[R]{\raisebox{1ex}[0pt][0pt]{\sffamily{\@author}}}
\makeatother

\begin{document}
% sezione (data)
\section{Lezione del 13-03-26}

% stili pagina
\thispagestyle{empty}
\pagestyle{fancy}

% testo
\subsection{Trasformata di Fourier continua}
La \textit{trasformata continua di Fourier} permette di portare un segnale $X(t)$ dal dominio del tempo $t$ nella sua rappresentazione sul dominio delle frequenze $X(F)$. In particolare:
\begin{definition}{Trasformata di Fourier continua}
Data una funzione $x(t)$, la sua trasformata di Fourier continua $X(f)$ è:
$$
X(f) = \int x(t) e^{-j 2 \pi f t} \, dt \quad \quad \text{(equazione di analisi)}
$$
\end{definition}
Chiamiamo la formula data anche \textit{equazione di analisi}, in quanto permette di passare allo spettro (dominio di frequenze) del segnale (\textit{analizzarlo}).

Notiamo che la trasformata vale per segnali $x(t)$ sia reali che complessi. Chiaramente, i segnali \textit{fisici} sono reali. La risultante $X(f)$, invece, è sempre complessa. 

Ad ogni segnale $x(t)$ corrisponde una e una sola rappresentazione in dominio delle frequenze $X(f)$, e vicevesa.

Per riportarci alla rappresentazione sul dominio del tempo $t$ originale, possiamo usare l'\textit{antitrasformata}:
\begin{definition}{Antitrasformata di Fourier continua}
Data una funzione $X(f)$, la sua antitrasformata di Fourier continua $x(t)$ è:
$$
x(t) = \int X(f) e^{j 2 \pi f t} \, df \quad \quad \text{(equazione di sintesi)}
$$
\end{definition}
Chiamiamo la formula data anche \textit{equazione di sintesi}, in quanto permette di ricostruire un segnale dal suo spettro (dominio di frequenze) (\textit{sintetizzarlo}).

\subsubsection{Spettri di ampiezza e fase}
Abbiamo detto che la trasformata $X(f)$ è sempre una funzione complessa. In particolare, può essere divisa nelle componenti polari modulo ed angolo:
$$
X(f) = |X(f)|e^{i \cdot \mathrm{arg}(X(f))}
$$
dove:
\begin{itemize}
	\item Il modulo $|X(f)|$ ci dà lo spettro di \textit{ampiezza} del segnale;
	\item L'argomento $\mathrm{arg}(X(f))$ ci dà lo spettro di \textit{fase} del segnale.
\end{itemize}

Per i segnali reali $\in \mathbb{R}$ vale la simmetria \textit{hermitiana} degli spettri:
$$
X(-f) = X^*(f)
$$
dove l'operatore asterisco rappresenta il coniugato. Questo si ha da:
$$
X^*(f) = \int \left( x(t) e^{-i 2 \pi f t} \, dt \right)*
= \int x(t)^* e^{i 2 \pi f t} \, dt
= \int x(t) e^{i 2 \pi f t} \, dt = X(-f)
$$
Visto che $x(t)^* = x(t)$ (se $x(t) \in \mathbb{R}$). Notiamo che questo non è soddisfatto automaticamente dai segnali $\in \mathbb{C}$ (in quanto non si può imporre quest'ultima proprietà). \hfill $\square$

Da questo deriva anche che si hanno particolari simmetrie dello spettro:
\begin{itemize}
	\item Lo spettro di \textit{ampiezza} del segnale ha simmetria \textit{pari}:
		$$
		|X(f)| = |X(-f)|
		$$
	\item Lo spettro di \textit{fase} del segnale ha simmetria \textit{dispari}:
		$$
		\mathrm{arg}(X(f)) = -\mathrm{arg}(X(-f))
		$$
\end{itemize}

\par\smallskip

Vediamo quindi alcuni esempi di \textit{spettri} di alcune funzioni importanti: 

\begin{itemize}
	\item \textbf{\textsf{Segnale esponenziale monolatero}} \\
Come primo esempio, vediamo lo spettro dell'\textit{esponenziale monolatero}, già vista in 3.2:
$$
x(t) = e^{-\frac{t}{\tau}} \cdot u(t)
$$
dove $u(t)$ è il comune gradino di Heaviside.

Questo si ottiene direttamente applicando la definizione di trasformata di Fourier:
$$
X(f) = \int_{-\infty}^{\infty} x(t) e^{-i 2 \pi f t} \, dt = \int_0^{\infty} e^{-\frac{t}{\tau}} e^{-i 2 \pi f t} \, dt = \int_0^\infty e^{-t\left( \frac{1}{\tau} + i 2 \pi f \right)} \, dt = \frac{e^{-t\left( \frac{1}{\tau} + i 2 \pi f \right)}}{\frac{1}{\tau} + i 2 \pi f} \Big|^{\infty}_0
$$
$$
= \frac{\tau}{1 + i 2 \pi f \tau}
$$

\item \textbf{\textsf{Impulso rettangolare}} \\
Come secondo esempio, vediamo lo spettro della funzione \textit{rettangolare}, anch'essa già vista in 3.2:
$$
x(t) = 
\begin{cases}
	1, \quad |t| < \frac{\tau}{2} \\
	\frac{1}{2}, \quad |t| = \frac{\tau}{2} \\
	0, \quad |t| > \frac{\tau}{2}
\end{cases}
$$

Questo si ottiene direttamente applicando la definizione di trasformata di Fourier:
$$
X(f) = \int_{-\infty}^{\infty} x(t) e^{-i 2 \pi ft} \, dt = \int_{-\frac{\tau}{2}}^{\frac{\tau}{2}} e^{-i 2 \pi f t} \, dt
= - \frac{1}{i 2 \pi f} e ^{-i 2 \pi f t} \Big|^{\frac{\tau}{2}}_{-\frac{\tau}{2}} = - \frac{1}{i 2 \pi f} \left( e^{-i \pi f \tau} - e^{i 2 \pi f \tau} \right)
$$
$$
= - \frac{1}{i 2 \pi f} \cdot - 2 i \sin (\pi f t) = \frac{\sin(\pi f \tau)}{\pi f} 
$$

La funzione che abbiamo incontrato prende il nome di \textit{seno cardinale}, indicato come:
$$
\mathrm{sinc}(\alpha) = \frac{\sin(\pi \alpha)}{\pi \alpha} \ \forall \alpha \neq 0, \quad \mathrm{sinc}(0) = 1
$$
dove in zero eliminiamo la discontinuità eliminabile.

Abbiamo quindi trovato la trasformata \textit{fondamentale}:
$$
x(t) = \mathrm{rect} \left( \frac{t}{\tau} \right) \, \Leftrightarrow \, X(f) = \frac{\sin(\pi f \tau)}{\pi f} \cdot \frac{\tau}{\tau} = \tau \mathrm{sinc}(f \tau)
$$
\end{itemize}

\subsubsection{Teorema di Parseval}
Riguardo alla trasformata di Fourier e alle definizioni di energia e potenza date in 3.2, possiamo dare l'importante teorema sull'\textit{energia}:
\begin{theorem}{}
L’energia di un segnale nel tempo è uguale all’energia in frequenza, ovvero:
$$
E_x = \int |x(t)|^2 \, dt = \int |X(f)|^2 \, df
$$
\end{theorem}

Questo si può dimostrare prendendo l'energia del segnale in dominio tempo:
$$
E_x = \int_{-\infty}^{\infty} x(t) x^*(t) \, dt
$$
e riportando la $x^*(t)$ nella forma di sintesi vista prima:
$$
= \int_{-\infty}^{\infty} x(t) \left( \int_{-\infty}^{\infty} X^*(f) e^{-i 2 \pi f t} \, df  \right) dt
= \int_{-\infty}^{\infty} X^*(f) \left( \int_{-\infty}^{\infty} x(t) e^{-i 2 \pi f t} \, dt \right) df  
$$
$$
= \int_{-\infty}^{\infty} X^*(f) X(f) \, df
$$

Chiamiamo quindi $\mathcal{E}_x(f)$ \textit{densità spettrale di energia}:
$$
\mathcal{E}_x(f) = |X(f)|^2
$$

\par\smallskip

Riguardo alla \textit{potenza}, vale invece:
$$
P_x = \lim_{T \rightarrow \infty} \frac{E_{xT}}{T} = \lim_{T \rightarrow \infty} \int \frac{|X(f)|^2}{T} \, df = \int \lim_{T \rightarrow \infty} \frac{|X(f)|^2}{T} \, df
$$
da dove troviamo la definizione di $\mathcal{P}_x(f)$ di \textit{densità spettrale di potenza}:
$$
\mathcal{P}_x(f) = \frac{|X(f)|^2}{T}  
$$

\end{document}
